\documentclass[12pt, twoside]{article}
\input{/Users/chris/github/course-files/docs/ib/preamble}

\fancyhead[LE]{\thepage}
\fancyhead[RO]{Name: \hspace{4cm} \,\\ \hspace{2.9cm} \,}
\fancyhead[LO]{La Scuola d'Italia / Dr.\ Huson / IB Math 12 HL \\* 9 September 2026}

\begin{document}
\subsubsection*{12.2 Classwork: Complex Numbers Restart (no calculator)}

Answer both problems in the spaces provided. Show your working. Give
all answers as exact values (surds and fractions of $\pi$), not decimals.

\begin{enumerate}[itemsep=0.15cm]

%--- Problem 1: Cartesian arithmetic, modulus, real unknowns [7 marks] ---
%--- Source: authored 2026-09-07 (HL syllabus 1.12) ---
\item[1.]

Let $z = 3 - 2i$ and $w = 1 + 4i$.

\begin{enumerate}[itemsep=0.15cm]
  \item Find $zw$, giving your answer in the form $a + bi$, where
  $a, b \in \mathbb{R}$.
  \item Find $\dfrac{z}{w}$, giving your answer in the form $a + bi$, where
  $a, b \in \mathbb{R}$.
  \item Find the exact value of $|z|$.
  \item Find the real numbers $x$ and $y$ such that
  $(x + yi)(2 - i) = 7 + 4i$.
\end{enumerate}

\begin{tikzpicture}
  \draw (-0.6,0) rectangle (11,14.5);
  \foreach \y in {1,2,...,13} \draw[dotted] (0,\y)--(10.5,\y);
\end{tikzpicture}

\newpage

%--- Problem 2: Non-real roots, conjugate pairs, sum and product [7 marks] ---
%--- Source: authored 2026-09-07 (HL syllabus 1.12) ---
\item[2.]

Consider the equation $z^2 - 4z + 13 = 0$, where $z \in \mathbb{C}$.

\begin{enumerate}[itemsep=0.15cm]
  \item Solve the equation, giving both roots in the form $a + bi$, where
  $a, b \in \mathbb{R}$.
  \item Explain why the two roots are complex conjugates of each other.

\end{enumerate}

The complex number $2 + i\sqrt{3}$ is one root of a \emph{different}
quadratic equation with real coefficients.

\begin{enumerate}[itemsep=0.15cm]
  \item[(c)] Without solving any equation, find the sum and the product of
  the two roots of that equation.
  \item[(d)] Hence write down a quadratic equation with real coefficients
  that has $2 + i\sqrt{3}$ as a root.
\end{enumerate}

\begin{tikzpicture}
  \draw (-0.6,0) rectangle (11,14.5);
  \foreach \y in {1,2,...,13} \draw[dotted] (0,\y)--(10.5,\y);
\end{tikzpicture}

\end{enumerate}
\end{document}